% ============================================================================
% Chapter 1 — Introduction
% ============================================================================
\chapter{Introduction}

The Caffeine Engine is a C++20 game engine built on top of SDL3. It is
designed around three principles: explicit control over memory, data-oriented
processing via an archetype-based ECS, and zero external runtime dependencies
beyond SDL3, ImGui, and ImNodes.

\section{How to Use This Document}

Each chapter in this reference targets one subsystem. Within each chapter
you will find:
\begin{itemize}[noitemsep]
  \item a short description of what the system does and when to use it;
  \item the relevant types, structs, and enums with their fields;
  \item the public API with parameter descriptions;
  \item complete, compilable code examples.
\end{itemize}

\section{Header Organisation}

All public headers live under \texttt{src/}. The main convenience header is:

\begin{lstlisting}
#include "Caffeine.hpp"   // includes everything
\end{lstlisting}

Or include individual modules:

\begin{lstlisting}
#include "core/Types.hpp"
#include "math/Vec2.hpp"
#include "ecs/World.hpp"
#include "ecs/Components.hpp"
#include "physics/PhysicsComponents2D.hpp"
#include "audio/AudioComponents.hpp"
#include "input/InputManager.hpp"
#include "scene/SceneComponents.hpp"
#include "script/CppScript.hpp"
\end{lstlisting}

\section{Namespaces}

\begin{longtable}{ll}
\toprule
\textbf{Namespace} & \textbf{Contents} \\
\midrule
\texttt{Caffeine}              & Types, math types \\
\texttt{Caffeine::ECS}         & World, Entity, all components \\
\texttt{Caffeine::Physics2D}   & Physics components and shapes \\
\texttt{Caffeine::Scene}       & Scene hierarchy components \\
\texttt{Caffeine::Audio}       & Audio components \\
\texttt{Caffeine::Input}       & InputManager, Key, Action, Axis \\
\texttt{Caffeine::Script}      & CppScript base class, registry \\
\bottomrule
\end{longtable}

\section{A Minimal Game Loop Sketch}

The following snippet illustrates the typical lifecycle of a game object:

\begin{lstlisting}
#include "ecs/World.hpp"
#include "ecs/Components.hpp"

using namespace Caffeine::ECS;

World world;

// create an entity
Entity player = world.create("Player");

// attach components
world.add<Transform>(player, Transform{
    .position = {100.0f, 200.0f, 0.0f},
    .rotation = {0.0f, 0.0f, 0.0f},
    .scale    = {1.0f, 1.0f, 1.0f}
});
world.add<Velocity2D>(player, Velocity2D{0.0f, 0.0f});

// per-frame update
ComponentQuery q;
q.require<Transform>().require<Velocity2D>();

world.forEach<Transform, Velocity2D>(q, [&](Entity e, Transform& t, Velocity2D& v) {
    t.position.x += v.x * dt;
    t.position.y += v.y * dt;
});
\end{lstlisting}

The chapters that follow explain every piece shown here in full detail.
